% Original Bengali editorial revision; unique theory integrated below.
\begin{topicbody}
\section{ভার্চুয়াল রিয়েলিটি}
\label{sec:ch1-vr}
\subsection{ভার্চুয়াল রিয়েলিটির ধারণা}
\definition{ভার্চুয়াল রিয়েলিটি}{কম্পিউটারের হার্ডওয়্যার ও সফটওয়্যার দিয়ে তৈরি এমন ত্রিমাত্রিক কৃত্রিম পরিবেশ, যেখানে ব্যবহারকারী বিশেষ যন্ত্রের সাহায্যে নিজেকে উপস্থিত মনে করে এবং পরিবেশের সঙ্গে কাজ করতে পারে, তাকে ভার্চুয়াল রিয়েলিটি (Virtual Reality বা VR) বলে।}
VR-এ কম্পিউটার গ্রাফিক্স, সেন্সর, শব্দ এবং প্রদর্শনযন্ত্র মিলিয়ে বাস্তবের মতো অভিজ্ঞতা তৈরি করা হয়। ব্যবহারকারী বাস্তবে এক জায়গায় থাকলেও অন্য কোনো পরিবেশে থাকার অনুভূতি পেতে পারে। মাথা ঘোরালে দৃশ্য বদলায়; নিয়ন্ত্রক নাড়ালে কোনো বস্তু ধরা বা সরানো যায়। সাধারণ ভিডিওতে দর্শক আগে ধারণ করা দৃশ্য দেখে; VR-এ ব্যবহারকারীর কাজ অনুযায়ী পরিবেশের প্রতিক্রিয়া বদলাতে পারে।
\chOneFigureExplain{\chaponeimg/vr_concept_scene.png}{ভার্চুয়াল পরিবেশের অভিজ্ঞতা}{fig:ch1-pair-16}{হেডসেট চোখের সামনে দৃশ্য দেখায়। সেন্সর ব্যবহারকারীর নড়াচড়া শনাক্ত করে এবং কম্পিউটার সেই অনুযায়ী দৃশ্য ও শব্দ বদলায়। নিয়ন্ত্রক দিয়ে পরিবেশের বস্তুর সঙ্গে কাজ করা যায়।}
\subsection{অভিজ্ঞতার মূল ভিত্তি}
VR-এর পরিচিত Three-I মডেল হলো Immersion, Interaction ও Imagination। এগুলোর সঙ্গে সম্পর্কিত আরেকটি গুরুত্বপূর্ণ ধারণা হলো Presence।
\begin{compactpoints}
\item \term{নিমগ্নতা (Immersion)}: দৃশ্য ও শব্দের প্রভাবে কৃত্রিম পরিবেশে ডুবে থাকার অভিজ্ঞতা। বিস্তৃত দৃশ্য, স্বাভাবিক শব্দ ও দ্রুত প্রতিক্রিয়া এই অভিজ্ঞতা বাড়ায়।
\item \term{মিথস্ক্রিয়া (Interaction)}: ব্যবহারকারীর কাজ অনুযায়ী পরিবেশের পরিবর্তন। যেমন, নিয়ন্ত্রক দিয়ে একটি ভার্চুয়াল বস্তু সরানো।
\item \term{কল্পনা (Imagination)}: বাস্তবে উপস্থিত নয় এমন পরিবেশ কল্পনা করে নির্মাণ করা।
\item \term{উপস্থিতির অনুভূতি (Presence)}: পরিবেশ কৃত্রিম জেনেও ব্যবহারকারীর নিজেকে সেখানে উপস্থিত মনে হওয়া। এটি Three-I-এর তৃতীয় ইংরেজি শব্দ নয়।
\end{compactpoints}
নিমগ্নতা পরিবেশের অভিজ্ঞতার দিকটি বোঝায়; মিথস্ক্রিয়া বোঝায় ব্যবহারকারীর কাজের প্রতিক্রিয়া। উপস্থিতি ব্যবহারকারীর মানসিক অনুভূতি। স্বাভাবিক দৃশ্য ও যথাযথ প্রতিক্রিয়া উপস্থিতির অনুভূতি জোরালো করতে পারে।
\chOneFigureExplain{\chaponeimg/vr_three_pillars.png}{Three-I ও উপস্থিতির অনুভূতি}{fig:ch1-pair-17}{পরিবেশ নির্মাণ, সেই পরিবেশে নিমগ্ন হওয়া এবং সেখানে কাজ করার সম্পর্ক চিত্রে দেখানো হয়েছে।}
\subsection{প্রধান উপাদান ও কাজ}
VR-এর যন্ত্রগুলো ইনপুট নেওয়া, তথ্য প্রক্রিয়াকরণ এবং অভিজ্ঞতা প্রদর্শনের কাজ করে। ব্যবস্থাভেদে যন্ত্রের সমন্বয় বদলায়; সব ব্যবস্থায় নিচের সব যন্ত্র একসঙ্গে লাগে না।
\begin{compacttable}
\begin{longtable}{|p{\dimexpr.30\linewidth-7pt\relax}|p{\dimexpr.70\linewidth-7pt\relax}|}
\hline\textbf{উপাদান} & \textbf{কাজ}\\\hline\endfirsthead
\hline\textbf{উপাদান} & \textbf{কাজ}\\\hline\endhead
হেড-মাউন্টেড ডিসপ্লে (Head-Mounted Display বা HMD) & চোখের সামনে দুই চোখের জন্য পৃথক দৃশ্য দেখিয়ে গভীরতার অনুভূতি তৈরি করে। অনেক হেডসেটে মাথার নড়াচড়া মাপার সেন্সরও থাকে।\\\hline
মোশন সেন্সর ও ট্র্যাকার & মাথা, হাত বা শরীরের অবস্থান, দিক ও নড়াচড়া মাপে।\\\hline
নিয়ন্ত্রক (Controller) & নির্দেশ দেওয়া, বস্তু নির্বাচন, ধরা ও সরানোর সুযোগ দেয়।\\\hline
ডেটা গ্লাভ (Data Glove) & হাত ও আঙুলের অবস্থান ও নড়াচড়া ইনপুট হিসেবে নেয়।\\\hline
বডি স্যুট (Body Suit) & সেন্সরের সাহায্যে শরীরের নড়াচড়া মাপে; কোনো নকশায় স্পর্শের প্রতিক্রিয়াও দিতে পারে।\\\hline
অডিও ব্যবস্থা & শব্দের দিক ও দূরত্ব বোঝাতে স্থানিক শব্দ (Spatial Audio) তৈরি করে।\\\hline
হ্যাপটিক যন্ত্র (Haptic Device) & কম্পন বা বলের প্রতিক্রিয়া দিয়ে স্পর্শ বা চাপের অনুভূতি দেয়।\\\hline
কম্পিউটার ও গ্রাফিক্স প্রসেসর (GPU) & ইনপুট বিশ্লেষণ করে দ্রুত দৃশ্য তৈরি ও হালনাগাদ করে।\\\hline
সিমুলেশন সফটওয়্যার & পরিবেশের বস্তু, গতি, পদার্থবিজ্ঞানের নিয়ম ও ব্যবহারকারীর কাজের প্রতিক্রিয়া নিয়ন্ত্রণ করে।\\\hline
\end{longtable}
\end{compacttable}
সেন্সর, নিয়ন্ত্রক, ডেটা গ্লাভ ও বডি স্যুট মূলত ইনপুট ও নড়াচড়ার তথ্য দেয়। কম্পিউটার, GPU ও সফটওয়্যার সেই তথ্য প্রক্রিয়াকরণ করে। HMD, স্পিকার ও হ্যাপটিক যন্ত্র ফল প্রদর্শন বা অনুভব করায়।
\remember{কিছু উদাহরণে দশটি সেন্সরযুক্ত বডি স্যুট দেখানো হয়। এটি নির্দিষ্ট নকশার তথ্য; সব বডি স্যুটে সেন্সরের সংখ্যা দশটি নয়।}
\chOneFigureExplain{\chaponeimg/vr_system_components.png}{VR ব্যবস্থার উপাদান}{fig:ch1-pair-18}{ইনপুট, প্রক্রিয়াকরণ ও অভিজ্ঞতা প্রদর্শনের যন্ত্রগুলো একসঙ্গে কাজ করে।}
\subsection{কাজের পদ্ধতি}
\begin{compactenum}
\item ব্যবহারকারীর মাথা, হাত বা শরীরের নড়াচড়া এবং নির্দেশ ইনপুট হিসেবে নেওয়া হয়।
\item ট্র্যাকার অবস্থান ও দিক নির্ণয় করে কম্পিউটারে পাঠায়।
\item সফটওয়্যার তথ্য বিশ্লেষণ করে পরিবেশের পরিবর্তন হিসাব করে।
\item ত্রিমাত্রিক গঠন থেকে দেখানোর উপযোগী দৃশ্য তৈরি হয়। এ কাজকে রেন্ডারিং (Rendering) বলে।
\item হেডসেটে দৃশ্য, স্পিকারে শব্দ এবং উপযুক্ত যন্ত্রে স্পর্শের প্রতিক্রিয়া পাওয়া যায়।
\item ব্যবহারকারীর পরবর্তী কাজ অনুযায়ী একই প্রক্রিয়ায় পরিবেশ আবার হালনাগাদ হয়।
\end{compactenum}
যেমন, মাথা ডানে ঘোরালে সেন্সর সেই পরিবর্তন জানায়। কম্পিউটার নতুন দিকের দৃশ্য দেখায়। ফলে ব্যবহারকারী চারদিকে তাকানোর অভিজ্ঞতা পায়।
\chOneFigureExplain{\chaponeimg/vr_working_process.png}{VR-এ ইনপুট থেকে প্রতিক্রিয়া}{fig:ch1-pair-19}{ব্যবহারকারীর কাজ থেকে তথ্য নিয়ে দৃশ্য তৈরি করা এবং পরবর্তী কাজ অনুযায়ী দৃশ্য বদলানোর চক্র দেখানো হয়েছে।}
VR নির্মাণের চারটি ধারণাগত অংশও বোঝা যায়: \term{Tracker} নড়াচড়া মাপে; \term{Simulator} পরিবেশের নিয়ম ও প্রতিক্রিয়া হিসাব করে; \term{Application} নির্দিষ্ট কাজ বা অভিজ্ঞতা নির্ধারণ করে; \term{Geometry} ত্রিমাত্রিক গঠন দেয়। এগুলো কাজের দিক; চারটি পৃথক যন্ত্র কেনার তালিকা নয়। CAD-এ তৈরি ত্রিমাত্রিক মডেল VR-এ ব্যবহার করা যায়; CAD-এর মূল আলোচনা অংশ~\ref{sec:ch1-production}-এ রয়েছে। মডেল তৈরি, ভার্চুয়াল পরিবেশে অনুশীলন এবং EHR-এ রোগীর তথ্য রাখা ভিন্ন কাজ।
\subtopichead{স্বাভাবিক অভিজ্ঞতার শর্ত}
\begin{compactpoints}
\item \term{কম বিলম্ব (Low Latency)}: নড়াচড়া ও দৃশ্য বদলানোর ব্যবধান কম হওয়া দরকার। বেশি বিলম্বে দেখা দৃশ্য ও শরীরের অনুভূতি না মিললে অস্বস্তি বা মাথা ঘোরা হতে পারে।
\item \term{উচ্চ ফ্রেম হার}: প্রতি সেকেন্ডে যথেষ্ট সংখ্যক দৃশ্য দেখালে চলাচল মসৃণ লাগে।
\item \term{দৃষ্টিক্ষেত্র (Field of View বা FOV)}: চোখের সামনে দৃশ্যের বিস্তার বেশি হলে নিমগ্নতা বাড়তে পারে।
\item \term{নির্ভুল ট্র্যাকিং}: মাথা ও হাতের নড়াচড়া সঠিকভাবে মাপলে কাজের প্রতিক্রিয়া স্বাভাবিক হয়।
\item শব্দের দিক ও দূরত্ব এবং উপযুক্ত স্পর্শের প্রতিক্রিয়া পরিবেশকে বাস্তবসম্মত করতে সাহায্য করে।
\end{compactpoints}
\subsection{সংক্ষিপ্ত ইতিহাস}
সিমুলেশন, কম্পিউটার গ্রাফিক্স, প্রদর্শনযন্ত্র ও সেন্সরের ধারাবাহিক উন্নয়ন থেকে আধুনিক VR গড়ে উঠেছে।
\begin{compacttable}
\begin{tabular}{|p{\dimexpr.20\linewidth-7pt\relax}|p{\dimexpr.80\linewidth-7pt\relax}|}
\hline\textbf{সময়} & \textbf{উন্নয়ন}\\\hline
1950s--1960s & Morton Heilig-এর Sensorama-তে দৃশ্য, শব্দ ও নড়াচড়া মিলিয়ে বহু ইন্দ্রিয়ের অভিজ্ঞতা তৈরির চেষ্টা করা হয়।\\\hline
1968 & Ivan Sutherland-এর প্রাথমিক HMD ব্যবস্থা মাথায় পরা প্রদর্শনযন্ত্রের বিকাশে গুরুত্বপূর্ণ ছিল।\\\hline
1980s & Jaron Lanier ও VPL Research-এর কাজে DataGlove এবং বাণিজ্যিক VR পরিচিতি পায়।\\\hline
1990s & CAVE-এ ঘরজুড়ে প্রক্ষেপিত দৃশ্য ব্যবহার করা হয়; গবেষণা ও VR arcade-এর চর্চা বাড়ে।\\\hline
2010s থেকে & Oculus Rift, HTC Vive, PlayStation VR ও Meta Quest-এর মতো যন্ত্র সাধারণ ব্যবহারকারীর কাছে VR পৌঁছে দেয়। স্বতন্ত্র হেডসেট নিজস্ব প্রসেসর দিয়ে কাজ করতে পারে।\\\hline
\end{tabular}
\end{compacttable}
\chOneFigureExplain{\chaponeimg/vr_history_timeline.png}{VR-এর বিকাশের ধারা}{fig:ch1-pair-20}{প্রাথমিক বহু ইন্দ্রিয়ভিত্তিক ব্যবস্থা থেকে আধুনিক হেডসেট পর্যন্ত বিকাশের ধাপ দেখানো হয়েছে।}
\subsection{সিমুলেশন, VR, AR ও MR}
\definition{সিমুলেশন}{বাস্তব কোনো ব্যবস্থার মডেল তৈরি করে তার আচরণ, ফলাফল বা কাজের পদ্ধতি পর্যবেক্ষণ ও অনুশীলন করাকে সিমুলেশন (Simulation) বলে।}
আবহাওয়ার মডেল, যানবাহনের চলাচলের মডেল কিংবা স্প্রেডশিটে হিসাবভিত্তিক মডেলও সিমুলেশন হতে পারে। সব সিমুলেশনে ব্যবহারকারী ভার্চুয়াল পরিবেশে উপস্থিত থাকার অভিজ্ঞতা পায় না।
\begin{compacttable}
\begin{tabular}{|p{\dimexpr.18\linewidth-7pt\relax}|p{\dimexpr.51\linewidth-7pt\relax}|p{\dimexpr.31\linewidth-7pt\relax}|}
\hline\textbf{ধরন} & \textbf{পরিবেশ ও কাজ} & \textbf{উদাহরণ}\\\hline
VR & কম্পিউটার-নির্মিত পরিবেশে উপস্থিতি ও মিথস্ক্রিয়ার অভিজ্ঞতা দেয়। & ভার্চুয়াল ল্যাব ও প্রশিক্ষণ\\\hline
AR (Augmented Reality) & বাস্তব দৃশ্যের ওপর ডিজিটাল বস্তু বা তথ্য যোগ করে। & ক্যামেরায় বাস্তব ঘরে আসবাবের মডেল দেখা\\\hline
MR (Mixed Reality) & বাস্তব ও ভার্চুয়াল বস্তু একই পরিসরে থেকে পারস্পরিক প্রতিক্রিয়ায় অংশ নিতে পারে। & মিশ্র বাস্তবতার প্রশিক্ষণ\\\hline
সাধারণ ভিডিও & আগে ধারণ করা দৃশ্য দেখায়; দর্শকের কাজ অনুযায়ী বস্তুর আচরণ বদলায় না। & প্রামাণ্যচিত্র বা ধারণ করা পাঠ\\\hline
\end{tabular}
\end{compacttable}
সাধারণ পর্দায় অভিজ্ঞতা দেখানো হলে তা নন-ইমারসিভ, বড় পর্দা বা প্রক্ষেপণভিত্তিক প্রশিক্ষণব্যবস্থায় সেমি-ইমারসিভ এবং সাধারণত HMD ও ট্র্যাকিং দিয়ে গভীর নিমগ্নতা তৈরি হলে ইমারসিভ VR বলা হয়। 360-degree ভিডিওতে চারদিকে তাকানো গেলেও সব ক্ষেত্রে বস্তু ধরা বা সরানো যায় না।
\term{মেটাভার্স (Metaverse)} হলো পারস্পরিক যোগাযোগের জন্য ভাগাভাগি করে ব্যবহার করা ভার্চুয়াল পরিবেশের ধারণা। কোনো একক অ্যাপ বা হেডসেটের নাম দিয়ে এর অর্থ সম্পূর্ণ হয় না।
\chOneFigureExplain{\chaponeimg/vr_vs_ar_simulation.png}{ভার্চুয়াল ও বাস্তব পরিবেশের সম্পর্ক}{fig:ch1-pair-21}{বাস্তব দৃশ্য বজায় রাখা, কৃত্রিম পরিবেশে প্রবেশ এবং মডেল চালিয়ে ফল দেখার পার্থক্য চিত্রে দেখানো হয়েছে।}
\subsection{ব্যবহারক্ষেত্র}
ঝুঁকিপূর্ণ, ব্যয়বহুল, দূরবর্তী বা জটিল পরিবেশে কাজ শেখার জন্য VR উপযোগী। তবে অনুশীলনের ফল বাস্তবে প্রয়োগ করতে প্রশিক্ষকের নির্দেশ এবং বাস্তব দক্ষতার মূল্যায়নও দরকার।
\chOneFigureExplain{\chaponeimg/vr_applications_map.png}{VR-এর ব্যবহারক্ষেত্র}{fig:ch1-pair-22}{শিক্ষা, চিকিৎসা, বিমান চালনা, প্রতিরক্ষা, শিল্প, স্থাপত্য, বিনোদন ও পুনর্বাসনে ব্যবহারের ক্ষেত্র দেখানো হয়েছে।}
\subtopichead{শিক্ষা}
জীববিজ্ঞানে হৃৎপিণ্ড, মস্তিষ্ক ও কঙ্কালের ত্রিমাত্রিক গঠন দেখা যায়। পদার্থবিজ্ঞানে তরঙ্গ, তড়িৎক্ষেত্র ও মহাকাশে গতি বোঝানো যায়। রসায়নে অণুর গঠন ও বিক্রিয়ার মডেল দেখা যায়। ভার্চুয়াল ল্যাবে বাস্তব রাসায়নিক দুর্ঘটনার ঝুঁকি ছাড়া অনুশীলন এবং প্রকৌশলে যন্ত্র খুলে জোড়া লাগানো শেখা যায়। নিজস্ব গতিতে বারবার অনুশীলন ও দূর থেকে অংশগ্রহণের সুযোগ থাকে।
\subtopichead{চিকিৎসা প্রশিক্ষণ ও সহায়তা}
শিক্ষানবিশ চিকিৎসক ভার্চুয়াল অপারেশন কক্ষে অস্ত্রোপচারের ধাপ, যন্ত্র চালানো এবং ল্যাপারোস্কোপিক পদ্ধতি অনুশীলন করতে পারে। অঙ্গ ও টিস্যুর মডেল দেখে গঠন বোঝে। ভুল ধাপ, সময়, চাপ ও নির্ভুলতার প্রতিক্রিয়া পেয়ে সংশোধনের সুযোগ পায়; অনুশীলনে বাস্তব রোগীর ক্ষতি হয় না।
পুনর্বাসনে নড়াচড়ার দক্ষতা ফিরে পেতে অনুশীলন, ব্যথা থেকে মনোযোগ সরানো এবং চিকিৎসকের তত্ত্বাবধানে ভীতি, উদ্বেগ বা PTSD-তে নিয়ন্ত্রিত অভিজ্ঞতা দেওয়ার কাজেও VR ব্যবহৃত হয়। চিকিৎসকদের যৌথ প্রশিক্ষণ ও চিকিৎসাসংক্রান্ত সিমুলেশনেও এর ব্যবহার আছে।
\chOneFigureExplain{\chaponeimg/vr_medical_training.png}{ভার্চুয়াল অস্ত্রোপচার প্রশিক্ষণ}{fig:ch1-pair-23}{ত্রিমাত্রিক অঙ্গের মডেল, যন্ত্র চালানো এবং কাজের প্রতিক্রিয়া শিক্ষানবিশকে অনুশীলনে সাহায্য করে।}
\subtopichead{বিমান চালনা, প্রতিরক্ষা ও জরুরি সেবা}
কৃত্রিম ককপিটে প্রশিক্ষণার্থী নিয়ন্ত্রণযন্ত্র ব্যবহার করে বিমান চালনা ও দিকনির্দেশ শেখে। ঝড়, কুয়াশা, বৃষ্টি, রাতের উড্ডয়ন, ইঞ্জিন বিকল হওয়া, রানওয়ের সমস্যা ও জরুরি অবতরণ বারবার অনুশীলন করা যায়। প্রশিক্ষক কাজ পর্যবেক্ষণ করে নম্বর ও প্রতিক্রিয়া দিতে পারেন। বাস্তব বিমান, অস্ত্র বা বিপজ্জনক পরিস্থিতি ব্যবহারের ব্যয় ও ঝুঁকি কমানো যায়। সামরিক কৌশল, উদ্ধারকাজ, অগ্নিনিরাপত্তা এবং জরুরি পরিস্থিতির সাড়াদানেও একই সুবিধা পাওয়া যায়।
\chOneFigureExplain{\chaponeimg/vr_flight_simulator_training.png}{বিমান চালনার অনুশীলন}{fig:ch1-pair-24}{আবহাওয়া ও জরুরি পরিস্থিতি বদলে প্রশিক্ষণ দেওয়া এবং কাজের মূল্যায়ন করা যায়।}
\subtopichead{ব্যবসা, স্থাপত্য ও শিল্প}
ভার্চুয়াল শোরুমে পণ্য দেখানো এবং ভবন নির্মাণের আগে ঘুরে দেখার অভিজ্ঞতা দেওয়া যায়। ক্রেতা ঘরের আকার, আলো, আসবাবের অবস্থান ও বারান্দার দৃশ্য দেখে সিদ্ধান্ত নিতে পারে। নকশার ত্রুটি নির্মাণের আগেই সংশোধন করা যায়। কারখানায় যন্ত্র চালনা, রক্ষণাবেক্ষণ, নিরাপত্তা মহড়া, নকশা পর্যালোচনা ও নমুনা পণ্যের পরীক্ষা করা যায়।
\subtopichead{বিনোদন, সংস্কৃতি ও যোগাযোগ}
গেম, অংশগ্রহণমূলক গল্প, ভার্চুয়াল কনসার্ট, জাদুঘর, ঐতিহাসিক স্থান ও পর্যটন অভিজ্ঞতায় VR ব্যবহৃত হয়। অবতার দিয়ে সামাজিক যোগাযোগ, সভা ও দূরবর্তী যৌথ কাজ করা যায়। ক্রীড়া প্রশিক্ষণ এবং সাংস্কৃতিক ঐতিহ্যের দৃশ্যায়নও এর ব্যবহার।
\subsection{সীমাবদ্ধতা ও নিরাপদ ব্যবহার}
ভালো হেডসেট, নিয়ন্ত্রক, সেন্সর ও কম্পিউটারের ব্যয় বেশি হতে পারে। পর্যাপ্ত যন্ত্র না থাকলে শিক্ষার্থী ও প্রতিষ্ঠানের সুযোগে বৈষম্য তৈরি হয়। দুর্বল সিমুলেশন ভুল ধারণা দিতে পারে; অতিরিক্ত গেম বা নিমগ্ন অভিজ্ঞতা আসক্তির ঝুঁকি বাড়াতে পারে।
দীর্ঘ সময় হেডসেট পরলে চোখে চাপ, ক্লান্তি, মাথাব্যথা, মাথা ঘোরা বা বমিভাব হতে পারে। দেখা নড়াচড়া ও শরীরের অনুভূতি না মেলা, বেশি বিলম্ব এবং কম দৃশ্য-হালনাগাদ হার অস্বস্তি বাড়াতে পারে। চারপাশের দেয়াল বা আসবাব না দেখে চললে আঘাত পাওয়ার আশঙ্কা থাকে।
VR ব্যবস্থা নড়াচড়া, অবস্থান, কণ্ঠস্বর, হাতের ভঙ্গি, চোখের গতিবিধি, ব্যবহারের ধরন ও অ্যাকাউন্টের তথ্য সংগ্রহ করতে পারে। এসব তথ্য অনিরাপদভাবে রাখা বা ভাগ করলে ব্যক্তিগত গোপনীয়তা ক্ষতিগ্রস্ত হয়।
\begin{compactpoints}
\item ব্যবহারের আগে চারপাশে খালি জায়গা নিশ্চিত করতে হবে; নিরাপদ সীমানা বা guardian ব্যবস্থা ব্যবহার করা যায়।
\item একটানা দীর্ঘ সময় ব্যবহার না করে বিরতি নিতে হবে। মাথা ঘোরা বা অস্বস্তি হলে হেডসেট খুলে বিশ্রাম নিতে হবে।
\item শিশুদের জন্য বয়সোপযোগী বিষয়বস্তু ও সময়সীমা মানতে হবে।
\item বিশ্বস্ত বিষয়বস্তু, নিরাপদ অ্যাকাউন্ট, অনুমতি নিয়ন্ত্রণ এবং তথ্য ব্যবহারের নীতি সম্পর্কে সচেতন থাকতে হবে।
\end{compactpoints}
\chOneFigureExplain{\chaponeimg/vr_health_safety.png}{VR ব্যবহারে স্বাস্থ্য ও নিরাপত্তা}{fig:ch1-pair-25}{ব্যবহারের সময়সীমা, নিরাপদ স্থান ও ব্যক্তিগত তথ্যের সুরক্ষা বিবেচনা করতে হয়।}
\subsection{অনুশীলন}

\qhead{MCQ}
\begin{multicols}{2}\begin{compactenum}
\item \begin{minipage}[t]{\linewidth}VR-এর পূর্ণরূপ কোনটি?\par \mcqoption{Virtual Reality}{Visual Router}{Video Record}{Virtual Radio}\chOneAnswerBox{ক}\par\end{minipage}

\item \begin{minipage}[t]{\linewidth}VR-এ ব্যবহারকারীর মাথার নড়াচড়া শনাক্ত করে কোন উপাদান?\par \mcqoption{Printer}{Motion sensor}{Scanner}{Modem}\chOneAnswerBox{খ}\par\end{minipage}

\item \begin{minipage}[t]{\linewidth}VR-এর গুরুত্বপূর্ণ বৈশিষ্ট্য কোনটি?\par \mcqoption{শুধু text দেখা}{কাগজ print করা}{Immersion ও Interaction}{offline storage}\chOneAnswerBox{গ}\par\end{minipage}

\item \begin{minipage}[t]{\linewidth}Head Mounted Display-এর কাজ কী?\par \mcqoption{paper scan করা}{network cable কাটা}{password তৈরি করা}{3D virtual scene দেখানো}\chOneAnswerBox{ঘ}\par\end{minipage}

\item \begin{minipage}[t]{\linewidth}Haptic device কী দেয়?\par \mcqoption{স্পর্শ/feedback-এর অনুভূতি}{শুধু print output}{শুধু text input}{শুধু storage}\chOneAnswerBox{ক}\par\end{minipage}

\item \begin{minipage}[t]{\linewidth}VR সবচেয়ে বেশি উপকারী কোন ক্ষেত্রে?\par \mcqoption{manual handwriting}{ঝুঁকিপূর্ণ training}{কাগজ ভাঁজ করা}{দেয়াল রং করা}\chOneAnswerBox{খ}\par\end{minipage}

\item \begin{minipage}[t]{\linewidth}VR ও সাধারণ video-এর মূল পার্থক্য কী?\par \mcqoption{video-তে interaction বেশি}{দুইটি একই}{VR-এ interaction থাকে}{VR কাগজে দেখা হয়}\chOneAnswerBox{গ}\par\end{minipage}

\item \begin{minipage}[t]{\linewidth}AR-এর কাজ কী?\par \mcqoption{সম্পূর্ণ virtual world তৈরি}{শুধু database রাখা}{sound বন্ধ করা}{বাস্তব environment-এর ওপর digital object বসানো}\chOneAnswerBox{ঘ}\par\end{minipage}

\item \begin{minipage}[t]{\linewidth}VR ব্যবহারে একটি স্বাস্থ্যঝুঁকি হলো—\par \mcqoption{motion sickness}{typing speed কমে যাওয়া}{paper jam}{printer heating}\chOneAnswerBox{ক}\par\end{minipage}

\item \begin{minipage}[t]{\linewidth}Surgical training-এ VR কেন উপকারী?\par \mcqoption{রোগীকে বাদ দিয়ে সব চিকিৎসা হয়}{রোগীর ঝুঁকি ছাড়াই practice করা যায়}{doctor দরকার হয় না}{medicine লাগে না}\chOneAnswerBox{খ}\par\end{minipage}

\item \begin{minipage}[t]{\linewidth}VR-এ প্রতিক্রিয়ার বিলম্ব বেশি হলে কী সমস্যা হতে পারে?\par \mcqoption{printer speed বাড়ে}{storage কমে যায়}{motion sickness হতে পারে}{network cable নষ্ট হয়}\chOneAnswerBox{গ}\par\end{minipage}

\item \begin{minipage}[t]{\linewidth}Data Glove কোন ধরনের তথ্য detect করতে সাহায্য করে?\par \mcqoption{কাগজের size}{বিদ্যুতের bill}{database password}{হাত ও আঙুলের movement}\chOneAnswerBox{ঘ}\par\end{minipage}

\item \begin{minipage}[t]{\linewidth}Presence বলতে কী বোঝায়?\par \mcqoption{virtual environment-এ নিজেকে উপস্থিত মনে হওয়া}{শুধু screen brightness}{file compression}{network address}\chOneAnswerBox{ক}\par\end{minipage}

\item \begin{minipage}[t]{\linewidth}VR-এ spatial audio কেন ব্যবহার করা হয়?\par \mcqoption{ফাইল print করতে}{শব্দের দিক ও দূরত্ব বোঝাতে}{keyboard lock করতে}{virus scan করতে}\chOneAnswerBox{খ}\par\end{minipage}

\item \begin{minipage}[t]{\linewidth}নিচের কোনটি VR-এর privacy risk-এর সঙ্গে সম্পর্কিত?\par \mcqoption{paper folding}{monitor cleaning}{movement ও usage data collection}{ink shortage}\chOneAnswerBox{গ}\par\end{minipage}
\end{compactenum}\end{multicols}

\qhead{সংক্ষিপ্ত প্রশ্ন}
\begin{compactenum}
\item ভার্চুয়াল রিয়েলিটি কী?
\item VR-এর তিনটি মূল ভিত্তি লেখ।
\item Head Mounted Display কী?
\item Data Glove-এর কাজ কী?
\item Haptic feedback কী?
\item VR system কীভাবে কাজ করে?
\item VR ও AR-এর মধ্যে পার্থক্য লেখ।
\item VR ও সাধারণ video-এর মধ্যে পার্থক্য লেখ।
\item চিকিৎসা প্রশিক্ষণে VR-এর ব্যবহার লেখ।
\item ফ্লাইট প্রশিক্ষণে VR কেন গুরুত্বপূর্ণ?
\item VR-এর তিনটি সুবিধা লেখ।
\item VR-এর তিনটি স্বাস্থ্যঝুঁকি লেখ।
\item VR ব্যবহারে কী কী সতর্কতা প্রয়োজন?
\end{compactenum}

\qhead{সৃজনশীল প্রশ্ন}
একটি training centre নতুন pilot ও surgeon-দের জন্য VR lab তৈরি করল। সেখানে headset, controller, motion sensor, data glove ও haptic device ব্যবহার করে virtual cockpit এবং virtual operation theatre তৈরি করা হলো। Trainee-রা ঝড়ের মধ্যে aircraft landing practice করে এবং surgical procedure বারবার অনুশীলন করে।

\begin{compactenum}
\item ক. VR কী?
\item খ. Haptic feedback বলতে কী বোঝায়?
\item গ. উদ্দীপকে ব্যবহৃত VR device-গুলোর কাজ ব্যাখ্যা কর।
\item ঘ. উদ্দীপকের আলোকে ঝুঁকিপূর্ণ training-এ VR-এর সুবিধা বিশ্লেষণ কর।
\end{compactenum}

\qhead{আরেকটি সৃজনশীল প্রশ্ন}
একটি স্কুলে science class-এর জন্য VR headset ব্যবহার করা হলো। শিক্ষার্থীরা human heart-এর 3D model দেখল, virtual lab-এ experiment করল এবং teacher-এর সঙ্গে virtual classroom-এ discussion করল। তবে কিছু শিক্ষার্থী দীর্ঘ সময় headset ব্যবহারের পর চোখে চাপ ও মাথা ঘোরার কথা জানাল।

\begin{compactenum}
\item ক. Immersion কী?
\item খ. VR কীভাবে learning experience উন্নত করে?
\item গ. উদ্দীপকে শিক্ষাক্ষেত্রে VR-এর ব্যবহার ব্যাখ্যা কর।
\item ঘ. উদ্দীপকের আলোকে VR ব্যবহারের সুবিধা ও সতর্কতা বিশ্লেষণ কর।
\end{compactenum}

\qhead{সৃজনশীল প্রশ্ন: ব্যবসা ও স্থাপত্য}
একটি real estate company buyer-দের জন্য VR walkthrough চালু করল। Buyer-রা headset পরে flat-এর room layout, lighting, furniture placement ও balcony view দেখে সিদ্ধান্ত নিতে পারছে। Company একই সঙ্গে buyer-এর viewing behaviour ও preference data সংরক্ষণ করছে।

\begin{compactenum}
\item ক. Head Mounted Display কী?
\item খ. VR ও AR-এর মূল পার্থক্য ব্যাখ্যা কর।
\item গ. উদ্দীপকে ব্যবসা ও স্থাপত্যে VR-এর ব্যবহার ব্যাখ্যা কর।
\item ঘ. উদ্দীপকের আলোকে VR ব্যবহারের সুবিধা এবং privacy risk বিশ্লেষণ কর।
\end{compactenum}

\end{topicbody}
